\newpage
\phantomsection
\label{prf:rational-multiplication-well-defined}

\begin{remark*}[Return]
\hyperref[lem:rational-multiplication-well-defined]{Return to Lemma}
\end{remark*}

\begin{lemma*}[Multiplication Is Well-Defined on Rational Classes]
Let $a,b,c,d,e,f,g,h\in\mathbb{Z}$ with $bdfh\neq 0$. If
\[
(a,b)\sim(e,f)
\qquad\text{and}\qquad
(c,d)\sim(g,h),
\]
then
\[
(ac,bd)\sim(eg,fh).
\]
\end{lemma*}

\begin{proof}
\textbf{Professional Standard Proof.}~\\
Assume
\[
(a,b)\sim(e,f)
\qquad\text{and}\qquad
(c,d)\sim(g,h).
\]
By the definition of $\sim$,
\[
af=be
\qquad\text{and}\qquad
ch=dg.
\]
To prove
\[
(ac,bd)\sim(eg,fh),
\]
it is enough to prove
\[
acfh=bdeg.
\]

Multiplying the equalities
\[
af=be
\qquad\text{and}\qquad
ch=dg
\]
gives
\[
(af)(ch)=(be)(dg).
\]
By associativity and commutativity of integer multiplication,
\[
acfh=bdeg.
\]
Therefore
\[
(ac,bd)\sim(eg,fh).
\]
Hence multiplication is well-defined on rational classes.
\end{proof}

\begin{proof}
\textbf{Detailed Learning Proof.}~\\

\textbf{Step 1.} Translate the hypotheses into cross-product equalities.

Assume
\[
(a,b)\sim(e,f)
\qquad\text{and}\qquad
(c,d)\sim(g,h).
\]
By the definition of the fraction-pair relation, these assumptions mean
\[
af=be
\qquad\text{and}\qquad
ch=dg.
\]

\textbf{Step 2.} Translate the desired conclusion into one equality.

The desired conclusion is
\[
(ac,bd)\sim(eg,fh).
\]
By the definition of $\sim$, this is equivalent to the cross-product equality
\[
(ac)(fh)=(bd)(eg).
\]
Equivalently, it is enough to prove
\[
acfh=bdeg.
\]

\textbf{Step 3.} Multiply the two hypothesis equalities.

From
\[
af=be
\qquad\text{and}\qquad
ch=dg,
\]
multiplying equal quantities gives
\[
(af)(ch)=(be)(dg).
\]

\textbf{Step 4.} Rearrange the products.

By associativity and commutativity of multiplication in $\mathbb{Z}$,
\[
(af)(ch)=acfh
\]
and
\[
(be)(dg)=bdeg.
\]
Therefore
\[
acfh=bdeg.
\]

\textbf{Step 5.} Return to the fraction-pair relation.

Since
\[
acfh=bdeg,
\]
we have
\[
(ac)(fh)=(bd)(eg).
\]
By the definition of $\sim$,
\[
(ac,bd)\sim(eg,fh).
\]
Thus multiplication is well-defined on rational classes.
\end{proof}

\begin{remark*}[Proof structure]
The proof is direct. The two hypotheses are converted into the
cross-product equalities defining equivalence of fraction pairs. These
equalities are multiplied together, and then associativity and commutativity
of integer multiplication rearrange the resulting products into the single
cross-product equality required for the products.
\end{remark*}

\begin{remark*}[Dependencies]~\\
\begin{itemize}
  \item \hyperref[def:fraction-pair-equality]{Definition~\ref*{def:fraction-pair-equality}}
  \item Basic arithmetic of the integers
  \item Associativity and commutativity of integer multiplication
\end{itemize}
\end{remark*}